%% LyX 2.4.4 created this file.  For more info, see https://www.lyx.org/.
%% Do not edit unless you really know what you are doing.
\documentclass[english]{beamer}
\usepackage{mathptmx}
\usepackage[T1]{fontenc}
\usepackage[latin9]{inputenc}
\usepackage{units}
\usepackage{amssymb}
\usepackage{graphicx}

\makeatletter
%%%%%%%%%%%%%%%%%%%%%%%%%%%%%% Textclass specific LaTeX commands.
% this default might be overridden by plain title style
\newcommand\makebeamertitle{\frame{\maketitle}}%
% (ERT) argument for the TOC
\AtBeginDocument{%
  \let\origtableofcontents=\tableofcontents
  \def\tableofcontents{\@ifnextchar[{\origtableofcontents}{\gobbletableofcontents}}
  \def\gobbletableofcontents#1{\origtableofcontents}
}

%%%%%%%%%%%%%%%%%%%%%%%%%%%%%% User specified LaTeX commands.
\usetheme{Warsaw}
% or ...

\setbeamercovered{transparent}
% or whatever (possibly just delete it)
\setbeamercovered{invisible} %makes overlays invisible


%gets rid of navigation symbols
\setbeamertemplate{navigation symbols}{}

\makeatother

\usepackage{babel}
\begin{document}
\title[Exam Review]{Exam Review: 9.1-9.5}
\subtitle{Calculus III | MTH 331~~|~~Section A~~|~~Fall 2012}
\author{Dr.\,James Ashe}
\titlegraphic{\includegraphics[height=1.5cm]{../../jcsu_bull.jpg}}
\institute{Johnson C. Smith University}
\date{{}}

\makebeamertitle

%\beamerdefaultoverlayspecification{<+->}
\begin{frame}{}

\textbf{Homework}

\textbf{9.1}: 21, 55, 60, 66, 68 

\textbf{9.2}: 4, 9, 18, 20, 29, 46, 48, 51-56, 58 

\textbf{9.3}: 12, 19, 23, 26, 47, 49 

\textbf{9.4}: 12, 14, 20, 24, 26, 29 

\textbf{9.5}: 9-14, 19-25, 27-28
\end{frame}

\begin{frame}{}

\begin{block}{Recall: Integration by Parts}

$\int u\,dv=uv-\int v\,du$
\end{block}
\begin{itemize}
\item <2->[ex.]Find $\int xe^{-x}\,dx$. \uncover<3->{Let $u=x$ so that,}\uncover<4->{
\[
\begin{array}{rcl}
u=x &  & dv=e^{-x}\,dx\\
du=x &  & v=-e^{-x}
\end{array}
\]
}\vspace{-10pt}
\item <5->[]$\int\underset{u}{\underbrace{x}}\cdot\underset{dv}{\underbrace{e^{-x}\,dx}}$\uncover<6->{$=\underset{u}{\underbrace{x}}\underset{v}{\underbrace{(-e^{-x})}}-\int\underset{v}{\underbrace{(-e^{-x})}}\,\underset{du}{\underbrace{dx}}$.}
\item <7->[]$=-xe^{-x}-e^{-x}+C$
\end{itemize}
\begin{block}<8->{$e^{x}$}
${\displaystyle \lim_{n\rightarrow\infty}\left(\frac{n+x}{n}\right)^{n}=e^{x}}$
\end{block}
\uncover<9->{In particular, ${\displaystyle \lim_{n\rightarrow\infty}\left(\frac{n}{n+1}\right)^{n}=e^{-1}}$.}
\begin{itemize}
\item <9->[]These facts come in useful on problems \textbf{9.4: }29, \textbf{9.5:
}12 and 14 (among others)
\end{itemize}
\end{frame}

\begin{frame}{Types of Problems}

\begin{block}{\textbf{9.1: }Overview of Sequences and Series}

Reading and writing the notation
\end{block}
\begin{itemize}
\item <2->Write: \textbf{ex. }Find an explicit formula for the sequence
$\{1,3,9,27,81,\dots\}$.
\item <3->Write: \textbf{ex. }Write $0.6+0.06+0.006+\dots$ in sigma notation.
\item <4->Read: \textbf{ex. }Find the first $4$ terms of the series ${\displaystyle \sum_{k=1}^{\infty}\frac{1}{2^{k}}}$. 
\end{itemize}
\begin{block}<4->{\textbf{9.2: }Sequences}

Determine the limit of a given sequence.
\end{block}
\begin{itemize}
\item <5->[ex.]Directly: \textbf{ex. }Find the limit of $\left\{ \frac{n^{3}}{n^{4}+1}\right\} $.
\item <6->[ex.]By comparing growth: \textbf{ex. }$\left\{ \frac{3^{n}}{n!}\right\} $;
$3^{n}\ll n!$.
\item <7->[ex.]By using the Squeeze Thm: \textbf{ex. }$?\leq\left\{ \frac{n\sin^{3}n}{n+1}\right\} \leq?$. 
\end{itemize}
\end{frame}

\begin{frame}{}

\begin{block}{\textbf{9.3: }Series}

Working with geometric and telescoping series.
\end{block}
\begin{itemize}
\item <2->Recognize and evaluate a geometric series: \textbf{ex. }${\displaystyle \sum_{k=0}^{\infty}\frac{3}{5^{k}}}$,
${\displaystyle \sum_{k=1}^{\infty}\left(\frac{9}{7}\right)^{k}}$.
\item <3->Recognize and evaluate a telescoping series: \textbf{ex. ${\displaystyle \sum_{k=1}^{\infty}\left(\frac{1}{k+2}-\frac{1}{k+3}\right)}$}
(no method of partial fractions).
\end{itemize}
\end{frame}

\begin{frame}{}

\begin{block}{\textbf{9.4: }Divergence and Integral Tests.}

Applying the divergence and integral tests.
\end{block}
\begin{itemize}
\item <2->Recognize the harmonic series, and recall that it diverges.
\item <3->Be able to manipulate series using basic properties: \textbf{ex.
}${\displaystyle \sum_{k=0}^{\infty}\left[2\left(\frac{3}{5}\right)^{k}+3\left(\frac{4}{9}\right)^{k}\right]}$.
\item <4->Use the Divergence Test: \textbf{ex. }${\displaystyle \sum_{k=2}^{\infty}\frac{1}{k\ln k}}$.
\item <5->Use the Integral Test: \textbf{ex. }${\displaystyle \sum_{k=1}^{\infty}\frac{1}{\sqrt[3]{k+10}}}$.
These problems may give you the $\frac{d}{dx}$ and $\int dx$. 
\end{itemize}
\end{frame}

\begin{frame}{}

\begin{block}{\textbf{9.5: }Ratio, Root, and Comparison Tests.}

Applying all these tests.
\end{block}
\begin{itemize}
\item <2->Ratio: Look at the ratio of $a_{k}$ and $a_{k+1}$ in a series.
\textbf{ex. }${\displaystyle \sum_{k=1}^{\infty}\frac{2^{k}}{k!}}$. 
\item <3->Root: Look at $\sqrt[n]{a_{k}}$ in a series. \textbf{ex. ${\displaystyle \sum_{k=1}^{\infty}\left(\frac{k+1}{2k}\right)^{k}}$}.
\item <4->Comparision: \textbf{ex. }${\displaystyle \sum_{k=1}^{\infty}\frac{\sin(\nicefrac{1}{k})}{k^{2}}}$,
\textrm{${\displaystyle \sum_{k=1}^{\infty}\frac{k^{2}}{\sin(\nicefrac{1}{k})}}$}
\end{itemize}
\uncover<5->{You will be provided with all major results}
\end{frame}

\end{document}
